\documentclass[10pt,UTF8]{book} %% ctexart

\title{\textbf{程序设计基础 (Java)}}
\author{钱锋\thanks{Email: strik0r.qf@gmail.com}${}^,$\thanks{
    西北工业大学软件学院, School of Software, Northwestern Polytechnical University, 西安 710072
}}

\usepackage{ctex}
\usepackage{graphicx}
\usepackage[toc]{multitoc}
\usepackage{booktabs}
\usepackage{longtable}
\usepackage{amsthm, amssymb, amsmath, mathrsfs, mhchem}
\usepackage{tikz}
\usetikzlibrary{decorations.markings, angles, quotes}
\usepackage{pgfplots}
\usepackage{tikz-3dplot}
\usepackage{extpfeil}
\usepackage{diagbox}
\usepackage{float}
\usepackage{hyperref}
\hypersetup{hidelinks,
    colorlinks = true,
    allcolors = black,
    pdfstartview = Fit,
    breaklinks = true}
\usepackage{caption}
\captionsetup[table]{labelsep=space} % 表
\captionsetup[figure]{labelsep=space} % 图
\usepackage{enumitem}
\usepackage{siunitx}

\usepackage{fancyhdr} % 用于自定义页眉页脚


% 设置页眉页脚样式
\fancypagestyle{plain}{%
    \fancyhf{} % 清空页眉页脚
    \fancyhead[RO,LE]{·\thepage·} % 页眉显示页码, RO表示奇数页右侧, LE表示偶数页左侧
    \fancyhead[LO]{\nouppercase{\rightmark}} % 页眉显示小节标题, LO表示奇数页左侧
    \fancyhead[RE]{\nouppercase{\leftmark}} % 页眉显示章节标题, RE表示偶数页右侧
    \renewcommand{\headrulewidth}{0.4pt} % 设置页眉横线的宽度
    \renewcommand{\footrulewidth}{0pt} % 取消页脚横线
}

\renewcommand{\headrule}{\hrule width\textwidth height\headrulewidth\vskip-\headrulewidth}

% % 取消奇偶页的页眉偏移
% \fancyhfoffset[RO,LE]{0pt}

% % 取消奇偶页的页眉偏移
% \fancyhfoffset[RO,LE]{0pt}

% 定义取消页眉的命令
\newcommand{\cancelheader}{%
    \fancyhead{} % 清空页眉
    \renewcommand{\headrulewidth}{0pt} % 取消页眉横线
    \renewcommand{\footrulewidth}{0pt} % 设置页脚横线的宽度
}

\renewcommand{\chaptermark}[1]{\markboth{第 \thechapter 章 \hspace{1em} #1}{}}
\renewcommand{\sectionmark}[1]{\markright{\thesection \, #1}}
\usepackage{titlesec} % 定义标题样式

% 设置 chapter 标题样式
\titleformat{\chapter}[hang]{\centering\heiti\Large\bfseries}{第\,\thechapter\,章}{1em}{}

% 定义 section 标题格式
\titleformat{\section}[hang]{\heiti\centering\large\bfseries}{\thesection}{1em}{}

% 定义 subsection 标题格式
\titleformat{\subsection}[hang]{\heiti\bfseries}{\textbf{\thesubsection}}{1em}{}

% 定义 subsubsection 标题格式
\setcounter{secnumdepth}{3}
\renewcommand\thesubsubsection{\arabic{subsubsection}.}
\titleformat{\subsubsection}[hang]{\kaishu}{\quad\quad\thesubsubsection\,\,}{0em}{}

% % 重新定义 textbf
% \let\oldtextbf\textbf
% \renewcommand{\textbf}[1]{{\heiti\oldtextbf{#1}}}

% % 在导言区重新定义 \normalsize 命令
% \makeatletter
% \renewcommand\normalsize{%
%    \@setfontsize\normalsize{10.5pt}{12pt}%
%    \abovedisplayskip 8\p@ \@plus2\p@ \@minus5\p@
%    \abovedisplayshortskip \z@ \@plus3\p@
%    \belowdisplayshortskip 6\p@ \@plus3\p@ \@minus3\p@
%    \belowdisplayskip \abovedisplayskip
%    \let\@listi\@listI}
% \makeatother



% 设置页边距和对齐
% \usepackage[
%     paperwidth=185mm,
%     paperheight=260mm,
%     top=35mm,
%     bottom=25mm,
%     left=18mm,
%     right=18mm,
%     footskip=15mm % 通过这里的值来调整页脚与正文内容的垂直距离
% ]{geometry}

\usepackage[
    paperwidth=210mm,
    paperheight=297mm,
    top=40mm,
    bottom=31.8mm,
    left=25.4mm,
    right=25.4mm,
    footskip=15mm % 通过这里的值来调整页脚与正文内容的垂直距离
]{geometry}

% \usepackage[
%     paperwidth=195mm,
%     paperheight=270mm,
%     top=40mm,
%     bottom=25mm,
%     left=23.5mm,
%     right=23.5mm,
%     footskip=15mm % 通过这里的值来调整页脚与正文内容的垂直距离
% ]{geometry}
\usepackage{mdframed}
\mdfsetup{
  linewidth=0.4pt,
  frametitlebackgroundcolor=white, % 或者 transparent
  frametitlefont=\heiti\bfseries,
  frametitleaboveskip=10pt,
  frametitlebelowskip=5pt,
  frametitlealignment=\raggedright % 新增此行
}
\usepackage{fontspec}
% 设置 Menlo 字体
\setmonofont{Menlo}
\usepackage{fancyvrb}
\usepackage{xcolor}
\usepackage{listings}

\definecolor{string}{HTML}{067D17}
\definecolor{comment}{HTML}{8C8C8C}
\definecolor{keyword}{HTML}{0033B3}
\definecolor{class_field}{HTML}{871094}

\lstset{breaklines}
%这条命令可以让LaTeX自动将长的代码行换行排版
\lstset{extendedchars=false}
%这一条命令可以解决代码跨页时，章节标题，页眉等汉字不显示的问题
\lstset{escapeinside={(*}{*)}}

\lstset{
    basicstyle=\small\ttfamily\heiti,
    numbers=left,
    numberstyle=\scriptsize\fontspec{Menlo}, % 使用 Menlo 字体
    stepnumber=1,
    numbersep=8pt,
    frame=leftline,
    xleftmargin=2em, % 调整代码块的左边界
    framexleftmargin=0pt, % 调整边框的位置
    breaklines=true,
    keywordstyle=\bfseries\color{keyword},          % keyword style
    commentstyle=\heiti\color{comment},       % comment style
    stringstyle=\color[HTML]{067D17},
    showstringspaces=false,
    % string literal style
    % escapeinside={\%*}{*)},            % if you want to add LaTeX within your code
    % morekeywords={}               % if you want to add more keywords to the set
}

\usepackage{smartdiagram}

\begin{document}

\newtheoremstyle{mytheoremstyle}
    {1.5ex}                                         % Space above
    {1.5ex}                                         % Space below
    {}                                              % Font for body
    {}                                              % Indent amount
    {\bfseries}                                     % Font for head
    {}                                              % Punctuation after head
    {0.5em plus 0.2em minus 0.1em}                  % Space after head
    {\thmname{#1}\thmnumber{ #2}.\thmnote{ (#3).}}

\theoremstyle{mytheoremstyle}
\newtheorem{definition}{定义}[section]
\newtheorem{example}{例}[section]
\newtheorem{exercise}{习题}[chapter]
\newtheorem{code}{程序清单}[section]
\newtheorem*{result}{运行结果}

\newtheoremstyle{my2theoremstyle}
    {1.5ex}                                         % Space above
    {1.5ex}                                         % Space below
    {\kaishu}                                              % Font for body
    {}                                              % Indent amount
    {\bfseries}                                     % Font for head
    {}                                              % Punctuation after head
    {0.5em plus 0.2em minus 0.1em}                  % Space after head
    {\thmname{#1}\thmnumber{ #2}.\thmnote{ (#3).}}

\theoremstyle{my2theoremstyle}
\newtheorem{thm}{定理}[section]
\newtheorem{law}{定律}[section]
\newtheorem{educt}{推论}
\newtheorem{prop}{命题}
\newtheorem{lemma}{引理}
\newtheorem{axiom}{公理}
\newtheorem{property}{性质}

\newtheoremstyle{my4theoremstyle}
    {1.5ex}                                         % Space above
    {1.5ex}                                         % Space below
    {}                                              % Font for body
    {}                                              % Indent amount
    {\bfseries}                                     % Font for head
    {}                                              % Punctuation after head
    {0.5em plus 0.2em minus 0.1em}                  % Space after head
    {\thmname{#1}.}

\theoremstyle{my4theoremstyle} \newtheorem*{sol}{解}

\newtheoremstyle{my3theoremstyle}
    {1.5ex}                                         % Space above
    {1.5ex}                                         % Space below
    {}                                              % Font for body
    {}                                              % Indent amount
    {\kaishu}                                       % Font for head
    {}                                              % Punctuation after head
    {0.5em plus 0.2em minus 0.1em}                  % Space after head
    {\thmname{#1}\thmnumber{ #2}.\thmnote{ (#3).}}

\theoremstyle{my3theoremstyle} \newtheorem*{remark}{注}
\newtheorem*{cmt}{评注}
\pagestyle{empty}
\begin{titlepage}
    \thispagestyle{empty}
    \centering
        \vspace*{2cm}
        \includegraphics[width=0.5\textwidth]{pic/npu_2.png}\par
        \vspace{1em}
        \includegraphics[width=0.5\textwidth]{pic/npu_1.png}\par
    \vspace{1em}
        \begin{center}
            \Huge \heiti \textbf{程序设计基础 (Java)}

            Fundamentals of Computer Programming
        \end{center}
        \vspace{17em}
        \begin{center}
        \songti
        \kaishu 软件学院 \, \heiti\textbf{钱锋} \quad \songti 编
        \vspace{0.5em}

    \today
    \end{center}
\end{titlepage}
\cleardoublepage
\maketitle
\cleardoublepage

\frontmatter
\newpage
\pagestyle{plain}
\makeatother

% 谨以此书, 敬伟大的数学家、伟大的数学教育家

% \begin{center}
%     Vladimir A. Zorich.
% \end{center}

% 这是一本抄来的数学书, 这本书的作者不是那个在海边发现贝壳的人, 只是一个沿着前人的足迹在展览馆里给贝壳拍照的后来人.

% 愿你在美妙的数学旅程中, 发现梦想, 追求真理和自由!


% \input{丛书前言.tex}


% \chapter*{第一版前言}
% \markboth{程序设计基础 (Java)}{第一版前言}
% \addcontentsline{toc}{chapter}{第一版前言}
% % 设置前言标题页的页码格式为empty, 即无页眉页脚
% \thispagestyle{empty}

% 本书是西北工业大学软件学院开设的《面向对象编程与设计》课程的先导参考资料,
% 旨在用尽可能短的时间帮助大家入门 Java 程序设计, 以更快地跟上《面向对象编程与设计》课程
% 直接以 Java 的面向对象特性为起点的教学进度.

% \begin{flushright}
%     \kaishu
%     钱锋

%     西北工业大学软件学院

%     2023 年 8 月
%     \songti
% \end{flushright}

\newpage
\thispagestyle{empty}

% 设置目录页的页码格式
\pagenumbering{roman} % 切换回罗马数字页码
\addtocontents{toc}{\protect\thispagestyle{empty}}
\pagestyle{plain}
{\tableofcontents}
\newpage
\thispagestyle{empty}
\cleardoublepage % 确保正文从奇数页开始


% 设置章节标题页的页眉和页脚为空白页样式
\makeatletter
\let\ps@plain\ps@empty
\makeatother

\mainmatter

\part{面向对象分析与设计}

\chapter{类与对象}

面向对象程序设计 (Object Oriented Programming, OOP)
就是使用对象进行程序设计.

对象是类的一个实例, 有状态和行为. 是属性 (数据、状态) 和处理属性的方法 
(行为、操作、服务) 的封装, 其中属性和方法作为一个不可分离的整体封装在一起, 
这也意味着方法是依赖于对象存在的, 没有游离于对象之外的方法.

类是一个创建对象的模板 (或某类对象的一个集合), 规定了该类中所有对象的属性和方法, 
对象是类的实例.

\chapter{Java 语言概述}

\quad\quad 关于 Java 语言的历史、计算机基础等内容我们在本章不做赘述, 我们直接从 Java 语言
的一些基本使用开始. 本章介绍一个简单的 Java 程序 \lstinline|Hello.java|, 并介绍 Java
的一些基础语法和 Java 语言的特点.

\section{Hello Java World!}


\begin{example}[Hello, Java!]
    \label{helloJavaWorld}
    这是一个经典的 hello 程序, 我们基于它来介绍 Java 程序的基本结构. 请编写一个
    \lstinline|Hello.java| 的 Java 类文件, 在 IDE 中输入以下代码, 然后尝试运行它:
    % \lstinputlisting[language=Java]{../IdeaProjects/Hello/src/Hello.java}
    这段程序的执行结果是向控制台输出信息 \lstinline|Hello Java World!|.
    接下来, 我们将从类的定义、方法的定义和 \lstinline|main| 方法、
    Java 语句和 \lstinline|.| 运算符三个角度来解释这段代码.
\end{example}

\subsection{类的定义}

Java 是纯面向对象的语言, 因此 Java 程序中的所有函数 
(我们在 Java 中习惯将其称为\textbf{方法}, method) 都必须放在\textbf{类} (class,
本书中常称为 Java class) 中来定义,
更彻底地说, Java 程序中的所有内容都必须在类中来组织.
类是 Java 程序的基本组成部分, 不管多大的 Java 工程, 都是由若干个 Java class 组织起来的.

在 Java 中使用 \lstinline|class|
关键字定义一个类, 具体的语法为
\begin{lstlisting}
(accessModifier) class className {

}
\end{lstlisting}
关于 Java class 的定义, 我们做几点说明:
\begin{enumerate}
    \item 在 \lstinline|class| 关键字前可以像例 \ref{helloJavaWorld} 一样添加 
    \lstinline|public| \textbf{访问修饰符} (access modifier), 我们会在后面详细介绍
    各种访问修饰符的作用及其使用方法, 现在你只需要知道, 访问修饰符用于控制程序的其他部分对
    这段代码的访问权限级别, Java 对访问权限的控制是封装性在 Java 中的一个具体实现.

    因此, \lstinline|public class Hello {}| 表示我们定义了一个
    访问权限为 \lstinline|public| 的类 \lstinline|Hello|, 而类中的所有元素的定义全部放在
    类名后面的大括号里.

    \item 在一个 Java 源文件中你可以定义多个类, 
    但是访问修饰符声明为 \lstinline|public| 的类 (这样的类称为\textbf{主类}) 最多
    只能有一个 (可以没有), 并且如果有的话, 源文件的文件名必须与主类的类名相同, 
    这就是我在前文中叫你把源文件命名为 \lstinline|Hello.java| 的原因.
    如果你没有注意到我的提示,
    把类 \lstinline|Hello| 声明在了其他源文件中, 那么你将会收到编译器的警告:
    \kaishu 错误: 类 \lstinline|Hello| 是公共的,
    应在名为 \lstinline|Hello.java| 的文件中声明. \songti

    一般来说, 我们约定一个源文件中仅编写一个类 (除非有测试需求).

    \item Java class 的命名规则为:
    \begin{enumerate}
        \item 类名必须以字母开头, 后面可以是字母和数字的任意组合, 一般用名词来给类命名;
        \item 采用 \lstinline|Pascal| 命名规则, 即每一个英文字母的首字母大写;
        \item 类名不能使用 Java 保留字;
        \item 大小写敏感.
    \end{enumerate}
\end{enumerate}

需要注意的是, 我们刚才仅仅展示了 Java 中对类的定义的语法中最简单的一个部分, 不过请放心,
作为一本 “喂饭式” 的参考资料, 当我们需要再次使用某个知识的时候, 我会贴心的给你复习一遍
(你甚至几乎不需要做任何额外的批注就能把以前学过的知识回顾一遍), 喂到你烦不胜烦. 因此
你将会一次, 又一次, 再一次的看到上面的介绍.

\subsection{方法的定义和 \lstinline|main| 方法}

\subsubsection{方法的定义}

在 Java 中定义一个方法的语法为
\begin{lstlisting}
(accessModifier) (static) returnType methodName(arg1Type arg1, ...) {

}
\end{lstlisting}
其中 \lstinline|accessModifier| 和 \lstinline|static| 是可选的, 
\lstinline|returnType| 指的是该方法的返回类型, 
\lstinline|methodName| 是方法名,
\lstinline|argType| 是该方法接受的
参数的类型, \lstinline|arg| 是该方法接收的参数名 (我们往往将其称为\textbf{形式参数}, 
你只要把它理解为函数 $f(x)$ 里的 $x$ 就行了, 其中, $x$ 是一个 \lstinline|argType| 类型
的变量). 方法中需要执行的指令序列称为方法体, 放在方法声明后的大括号中.

\subsubsection{\lstinline|main| 方法的定义}

在 \lstinline|main| 方法的定义中, 我们使用了关键字 \lstinline|public| 
和 \lstinline|static|, 并且指定 \lstinline|main| 方法没有返回值 (\lstinline|void|).
现在你不需要去管形式参数 \lstinline|String[] args| 是什么 (老实说, 直到我在写这一章的时候
我依然不知道它是什么), 我们会在后面详细介绍.
你需要知道的是:
\begin{enumerate}
    \item \lstinline|main| 方法是 Java 应用程序的执行入口, 它的固定声明格式为
    \begin{lstlisting}
class className {
    public static void main(String[] args) {}
}
    \end{lstlisting}
    现阶段, 请你向八股文一样把它死背住, 随着我们学习的深入, 你会逐渐理解这里面每一个关键词的含义.
    
    \item 每一个类都可以有 \lstinline|main| 方法,
    即 \lstinline|main| 方法也可以写在非 \lstinline|public| 类中, 然后指定运行非
    \lstinline|public| 类, 那么此时程序的入口就将是该非 \lstinline|public| 
    类的 \lstinline|main| 方法.
    
    \item 需要注意的是, \lstinline|main| 方法的拼写是 main, 许多人会在不经意间错拼成 mian,
    因此当找了很久找不到错误时, 不妨看看你有没有犯这一个低级的拼写错误.
\end{enumerate}

% 对象名、变量名和方法名采用 \lstinline|camelCase| 命名规则, 即第一个单词全小写, 
% 第二个单词开始每个单词的首字母大写. 除此之外, 类名和方法名不能使用 Java 的关键字, 
% 也不能使用数字开头,
% 类名和方法名是大小写敏感的, 也就是说, \lstinline|method1| 和 \lstinline|Method1|
% 是两个不一样的方法.

\subsection{Java 语句和 \lstinline|.| 运算符}

\lstinline|System.out.println("Hello Java World!");| 是 \lstinline|main|
方法中唯一的指令, 它的作用是向控制台输出字符串 \lstinline|Hello Java World|,
在这里你不需要掌握太多, 需要你明确的是:

\begin{enumerate}
    \item 一个 Java 语句需要用 \lstinline|;| 结尾, 如果不加分号的话是会报错的,
    语句末尾缺少分号或者输入全角分号 \lstinline|: | 是初学者常犯的错误.
    
    \item Java 中利用大括号 \lstinline|{}| 来划分代码的结构, 先前提到的类体、方法体都是写在类声明、
    方法声明后的大括号中的, Java 中的大括号总是成对出现的, 请你养成成对输入大括号的习惯.
    关于大括号的次行风格和行尾风格我们不再单独介绍,
    我本人采用的是行尾风格 (听不懂吗?
    没有关系, 如果你不知道我在说什么的话, 你就按照我的示例代码来输入大括号就可以了).
    
    \item Java 是面向对象的语言, 因此我们通常用 \lstinline|object.method()| 来调用
    对象的方法, 用 \lstinline|object.attribute| 来访问对象的属性.
\end{enumerate}

\section{Java 的基础语法}

\subsection{Java 转义字符}

\begin{table}[H]
    \caption{Java 中的转义字符}
    \centering
    \begin{tabular}{l l p{6cm}}
        \hline
        转义字符 & 用途 & 备注 \\
        \hline
        \lstinline|\t| & 制表位 & 用于实现对齐 \\
        \lstinline|\n| & 换行符 & 将光标移动到下一行的开头 \\
        \lstinline|\"| & 双引号 \\
        \lstinline|\'| & 单引号 \\
        \lstinline|\r| & 回车   & 将光标移动到当前行的开头, 新输出的内容会顶掉原先输出的内容 \\
        \hline
    \end{tabular}
\end{table}

\subsection{Java 注释}

用于注解、说明、解释程序的文字就是注释 (comment), 
被注释的文字不会被 \textbf{Java 虚拟机} (Java Virtual Machine, JVM) 解释执行,
因此你可以放心地往代码中添加注释,
注释能有效地提高代码的可读性并
降低团队成员间的沟通成本, 养成良好的注释习惯是很重要的. 不写注释的代码就像从头到尾只有
逻辑符号和运算过程, 而没有任何文字说明的数学文献一样粗暴.

良好的编程习惯是\textbf{自顶向下的实现}, 你首先需要做一个顶层的设计和规划, 然后再编码阶段
先通过注释写出你的实现策略和实现思路, 然后再逐步用代码去实现你的想法. 你的实现思路可能是由
若干个子模块组成的, 对于这些子模块, 采用同样的自顶向下实现方法, 直到它们被分解地足够基本
以至于每一个模块只实现一个很简单的功能为止.

在 Java 中实现注释的方式有三种: 单行注释、多行注释和 javadoc 文档注释, 它们的使用方法如下:
\begin{lstlisting}
// Single line comment

/*
  Multi line comments
  System.out.println("This line won't be executed.");
  */

/**
 * javadoc comments
 * @javadocTag this comment will be compliered by javadoc tool.
 */
\end{lstlisting}
关于注释的使用, 我们说明以下几点:
\begin{enumerate}
    \item 每一个 \lstinline|*/| 都会被认为是多行注释的结束, 因此
    多行注释内不可嵌套多行注释, 所以不要简单地把一段代码用 \lstinline|/* */| 括起来
    就认为它们已经被全部注释掉了.
    \item 文档注释的注释内容可以被 JDK 提供的 javadoc 工具所解析, 生成一套以网页文件的形式
    体现的程序说明文档, 其中 \lstinline|@javadocTag| 是 javadoc 标签, 常见的 javadoc
    标签见下表所示:
    \begin{table}[H]
        \centering
        \caption{常用 Javadoc 标签}
        \begin{tabular}{p{2.5cm} p{5cm} p{7cm}}
            \hline
            \textbf{标签} & \textbf{描述} & \textbf{示例} \\
            \hline
            \lstinline|@param| & 用于描述方法的参数, 提供参数的名称和描述.  & 
            \begin{lstlisting}
/**
 * @param a 第一个加数
 * @param b 第二个加数
 */
            \end{lstlisting} \\
                              \hline
            \lstinline|@return| & 用于描述方法的返回值类型和意义.  & 
            \begin{lstlisting}
/**
 * @return 学生的姓名
 */
            \end{lstlisting} \\
                              \hline
            \lstinline|@throws| & 用于描述方法可能抛出的异常类型和原因.  & 
            \begin{lstlisting}
/**
 * @throws Exception
 */
            \end{lstlisting} \\
                              \hline
            \lstinline|@deprecated| & 标记方法或类已经过时, 提醒开发者不再使用.  & 
            \begin{lstlisting}
/**
 * @deprecated 使用新方法代替
 */
            \end{lstlisting} \\
                                    \hline
            \lstinline|@see| & 引用其他类、方法、字段或文档.  & 
            \begin{lstlisting}
/**
 * 查看其他类的文档
 * @see OtherClass
 */
            \end{lstlisting} \\
                           \hline
        \end{tabular}
        \label{tab:javadoc-tags}
    \end{table}
    \item 类、方法的注释, 要用 javadoc 文档注释来写, 非 javadoc 文档注释内容,
    往往是为了便于代码的维护, 即告诉读者这段代码的思路以及相关的注意事项.
\end{enumerate}

\section{Java 语言的特点}

Java 的设计者编写了一个白皮书\footnote{可以在 
\url{www.oracle.com/technetwork/java/langenv-140151/html} 上找到.}来解释设计 
Java 的初衷以及完成这些目标的情况. 接下来我们介绍其中的四个关键术语, 
分别是面向对象、健壮性、可移植性和解释性, 它们是现阶段你需要掌握的 Java 的四个核心特点.

\subsection{面向对象}

Java 是一种纯粹的面向对象编程语言, 这意味着一切皆为对象. 对象是程序的基本构建块, 具有状态和行为. 
相比于传统的面向过程编程, 面向对象编程使得代码更模块化, 易于维护和扩展. 在 Java 中, 
类和对象的概念是核心, 而且通过接口的灵活运用, 实现了对多继承的替代方案. 

\subsection{健壮性}

Java 的健壮性体现在多个方面. 首先, Java 强调编译期和运行时的错误检测, 
通过严格的语法检查和类型检查, 能够在程序执行前发现潜在问题. 其次, 
Java 的垃圾回收机制有助于防止内存泄漏, 大大提高了程序的稳定性. 此外, 
异常处理机制也使得开发者能够更好地处理错误情况, 确保程序在面对异常时依然能够稳健运行. 

\subsection{可移植性}

Java 的可移植性是通过字节码和 Java 虚拟机 (JVM) 实现的. Java 程序在编译时生成字节码, 
而不是直接生成机器码. 这种字节码可以在任何装有相应 Java 虚拟机的平台上运行. 
这使得 Java 程序无需重新编写即可在不同的操作系统和硬件上执行, 为跨平台开发提供了强大的支持. 

\subsection{解释性}

Java 是一种解释性语言, 它的程序在运行时通过 Java 虚拟机进行解释执行. 
这种特性为 Java 带来了很多优势. 首先, Java 程序可以逐行执行, 
无需预先编译成本地机器码, 使得开发和调试更为灵活. 
其次, 解释性也有助于实现平台无关性, 因为相同的字节码可以在不同平台上被解释执行, 无需修改源代码. 

\vspace*{1cm}

除此之外,
Java 在设计上注重安全性. 通过强大的安全管理器, Java 可以在运行时防止恶意代码的执行. 
Applet 的沙箱机制是其中的一例, 它限制了从网络上下载的代码对本地系统的访问权限, 
防止了潜在的安全威胁. 
Java 还拥有庞大而丰富的生态系统, 有着丰富的标准库和框架. 
这使得开发者可以轻松地利用现有的工具和资源来加速开发过程. 从企业级应用到移动应用, 
Java 生态系统涵盖了各种领域, 为开发者提供了广泛的选择和支持. 

\newpage
\thispagestyle{empty}

\part{Java 程序设计}

\chapter{变量与运算符}

\quad\quad 变量是赋给值的标签, 指向了内存中数据的存储单元. Java 也用变量来存储值.

\section{变量}
\section{数值数据类型及其操作}

\newpage
\section{数值数据类型变量的类型转换}

两个不同数值数据类型的操作数进行二元运算的时候, 或者把一个类型的变量赋值给另一个类型的变量
的时候, 就需要进行类型转换. 在 Java 中, 数值数据类型类型转换分为自动类型转换和强
制类型转换. \textbf{类型转换}是将一种数据类型的值转换成另一种数据类型的值的操作:
\begin{enumerate}[label=(\arabic*), itemsep=0pt]
    \item 从占用字节数较小的类型转换成占用字节数较大的类型的转换称为\textbf{扩大类型}
    (widening a type), 在 Java 中, 扩大类型一般是自动完成的;
    \item 从占用字节数较大的类型转换成占用字节数较小的类型的转换称为\textbf{缩小类型}
    (narrowing a type), 在 Java 中如果想要缩小类型, 就必须显式地进行类型转换.
\end{enumerate}

\subsection{自动类型转换}

在 Java 中, 数值数据类型之间的自动转换遵循\textbf{自动提升原则}, 
即在一个表达式中有多个数参与运算时, 先把这些参与运算的数转换为容量最大的类型, 
再计算表达式的结果. 怎么理解呢? 大烧杯能装的东西肯定比小烧杯能装的东西要多,
因此占用字节数更大的数据类型拥有更多的二进制位, 也就有更大的数据表示范围,
所以我们在计算时先把所有的数据都 “倒进最大容量规格的那一个烧杯中”.

\subsection{强制类型转换}

% \section{Java}

% \subsection{Java 的基本特性}
% \subsubsection{面向对象}
% \subsubsection{健壮性}
% \subsubsection{可移植性}
% \subsubsection{解释性}

% \part{Java 的面向对象特性}

% \chapter{Java 程序设计语言}

% Java 是一门面向对象的程序设计语言, 一个 Java 程序可以认为是一系列对象的集合, 这些对象通过调用彼此的方法来协同工作.

% \section{Java 基础语法}

% \begin{example}
%     我们从经典的 Hello World 程序开始我们的 Java 旅途, 并以它作为示例来介绍 Java 程序的结构和框架.
%     \begin{lstlisting}[language=java]
% // HelloWorld.java
% package com.example.helloworld;  
    
% public class HelloWorld {
%     public static void main(String[] args) {
%         System.out.println("Hello World!");  
%     }  
% }
%     \end{lstlisting}
% \end{example}


% 在上面的示例中, public 是访问修饰符, 用于控制程序的其他部分对这段代码的访问级别;
% static 关键字 
% void 返回类型说明 
% main 是方法名, Java 虚拟机总是从指定类中的 main 方法开始执行, main 方法的访问修饰符必须声名为 public.

% 关于 Java 的基础语法, 我们做几点说明:

% \begin{itemize}
%     \item 大小写敏感;
%     \item 类名应该采用 CamelCase 命名法, 即各个单词的首字母大写, 这种命名法的英文名很容易理解 —— 那些大写的字母就像驼峰一样凸起;
%     \item 方法名应该以小写字母开头, 第二个单词开始每个单词的首字母大写;
%     \item Java 源文件名必须和类名相同, 文件的后缀名为 .java, 因此 HelloWorld 的文件命名应为 HelloWorld.java;
%     \item 所有的 Java 程序都从 main 方法开始执行.
% \end{itemize}

% \chapter{封装性在 Java 的编程实现}

% % \section{Java Class}
% % \subsection{创建一个新的数据类型: class}

% % 在 Java 中, 一个类通过 class 关键字进行定义, 类的定义包括类声明和类体两个部分.

% % \subsubsection{类声明: 创建一个新的对象类型}

% % 在 Java 中声明一个类的语法如下:
% % \begin{lstlisting}[language=Java]
% % public class ClassName {
% %     private int data1;                      //成员变量声明
% %     ....
% %     public int method1(int arg1,....){
% %         ....
% %     }
% %     ....
% % }
% % \end{lstlisting}

% \section{使用预定义类}

% \subsection{对象与对象变量}

% 对象变量并不实际包含一个对象, 它只是引用了一个对象.

% \subsection{Java 类库中的 LocateDate 类}

% 事实上, 时间和度量和时间节点的命名是一件非常复杂的工作.

% \subsection{更改器方法与访问器方法}

% 访问并修改对象的方法称为\textbf{更改器方法} (mutater method), 只访问对象而不修改对象的方法称为\textbf{访问器方法} (accessor method).


% \begin{example}[当前月的日历]
%     编写程序显示当前月的日历, 如
%     \begin{lstlisting}
% Mon Tue Wed Thu Fri Sat Sun
%                           1 
%   2   3   4   5   6   7   8 
%   9  10* 11  12  13  14  15 
%  16  17  18  19  20  21  22 
%  23  24  25  26  27  28  29 
%  30  31
%     \end{lstlisting}
%     其中, “今天” 处要加星号.
%     \begin{lstlisting}[language=java]
% import java.time.*;

% public class CalendarTest
% {
%     public static void main(String[] args){
%         // 获取当前日期
%         LocalDate date = LocalDate.now();        
%         // 获取当前月份
%         int month = date.getMonthValue();        
%         // 获取当前日期的日
%         int today = date.getDayOfMonth();        

%         // 将日期设置为本月的第一天
%         date = date.minusDays(today - 1);        
%         // 获取本月第一天是星期几
%         DayOfWeek weekday = date.getDayOfWeek(); 
%         int value = weekday.getValue();          

%         // 打印星期标题
%         System.out.println("Mon Tue Wed Thu Fri Sat Sun");  
%         // 打印月份第一天之前的空格
%         for (int i=1; i < value; i++)
%             System.out.print("    ");            
        
%         // 循环打印月份的每一天
%         while (date.getMonthValue() == month) {  
%             // 使用printf格式化输出日期, 保持3位宽度
%             System.out.printf("%3d", date.getDayOfMonth()); 
%             // 如果日期是今天, 则在日期后面添加"*"
%             if (date.getDayOfMonth() == today)
%                 System.out.print("*");           
%             else
%                 System.out.print(" ");
%             // 将日期向后推移一天
%             date = date.plusDays(1);             
%             // 如果日期是星期一, 换行
%             if (date.getDayOfWeek().getValue() == 1)
%                 System.out.println();            
%         }
        
%         // 如果最后一天不是星期一, 换行
%         if (date.getDayOfWeek().getValue() != 1)
%             System.out.println();                   
%     }
% }        
%     \end{lstlisting}
% \end{example}

% \chapter{接口}

% \subsection{继承父类与实现接口的区别}

% \begin{lstlisting}
% package InterfaceLearning.hspInterface;

% public class ExtendsVsImplements {
%     public static void main(String[] args) {
%         LittleMonkey wuKong = new LittleMonkey("悟空");
%         wuKong.climbing();
%         wuKong.swimming();
%     }
% }

% class Monkey {
%     private String name;
%     public Monkey(String name) {
%         this.name = name;
%     }
%     public String getName() {
%         return this.name;
%     }
%     public void climbing() {
%         System.out.println(name+"会爬树....");
%     }
% }

% class LittleMonkey extends Monkey implements Fishable {
%     public LittleMonkey(String name) {
%         super(name);
%     }

%     @Override
%     public void swimming() {
%         System.out.println(this.getName()+"学会了游泳....");
%     }
% }

% interface Fishable {
%     void swimming();
% }
% \end{lstlisting}

% \begin{enumerate}
%     \item 当子类继承了父类, 就自动拥有了父类的功能
%     \item 如果子类需要扩展功能, 可以通过实现接口的方式扩展
%     \item 实现接口是对 Java 单继承机制的补充
%     \item 继承的主要价值在于: 解决代码的复用性和可维护性
%     \item 接口的价值在于: 设计好各种规范, 让其他类去实现这些方法, 使代码更加灵活
%     \item 接口比继承更加灵活
%     \item 继承 is-a
%     \item 接口 like-a
%     \item 接口在一定程度上实现代码解耦 (接口规范性+动态绑定)
% \end{enumerate}

% \part{Java 高级编程}

% \chapter{Java 多线程编程}

% \section{Java 线程}
% \subsection{中央处理器 (CPU) }

% \subsubsection{CPU 的组成}

% \begin{enumerate}[label=(\arabic*)]
%     \item \textbf{运算器}: 执行逻辑运算, 对数据进行加工处理: 
%     \item \textbf{控制器}: 负责协调并控制计算机各个部件执行程序的指令序列, 
%     包括取指令、分析指令和执行指令: 
%     \item \textbf{寄存器}: 保存运算或指令的一些临时数据: 
% \end{enumerate}

% \subsubsection{单核与多核 CPU}

% \begin{enumerate}[label=(\arabic*)]
%     \item \textbf{单核}: 在一个时间单元内只能执行一个任务, 同时间段内有多个任务需要 CPU
%     去执行时, CPU 通过交替去执行多个任务中的一个任务来实现, 但是由于其执行速度快, 
%     用户很难觉察: 
%     \item \textbf{多核}: 多核的 CPU 能同时执行多个任务. 
% \end{enumerate}

% \subsection{程序与进程}

% \subsubsection{程序}

% 一组计算机能够识别和执行的指令, 实现特定任务的运行. 

% 开发者所编写的代码称之为程序, 程序就是按照算法和数据结构构件的代码, 一组
% 数据和指令集, 是一个静态的概念. 

% \subsubsection{进程}

% \begin{definition}[进程]
%     应用程序在执行过程中占用的独立的系统资源 (内存、CPU、磁盘、网络) 称为\textbf{进程} (Process) . 
% \end{definition}

% \begin{enumerate}[label=(\arabic*)]
%     \item CPU 从硬盘中读取一段程序到内存中, 该执行程序的实例就叫做进程: 
%     \item 一个程序如果被 CPU 多次读取到内存中, 变成多个独立的进程, 将程序运行起来, 
%     我们称之为进程: 
%     \item 进程是程序的一次执行过程, 是一个动态的概念: 
%     \item 进程存在生命周期, 进程随着程序运行的终止而销毁. 
% \end{enumerate}

% \subsubsection{线程}

% \begin{definition}[线程]
%     运行于进程中的任务单元称为\textbf{线程} (Thread) , 是进程中实际运作的单位. 
% \end{definition}
% \begin{enumerate}[label=(\arabic*)]
%     \item 通常在一个进程中可以包含若干个线程: 
%     \item 线程是 CPU 调度和执行的最小单位: 
%     \item 一个进程可以有多个线程, 例如, 在 Bilibili 弹幕视频网看视频时可以同时看图像、听声音、看弹幕等: 
%     \item 很多线程都是模拟出来的, 真正的多线程是指有多个 CPU, 或多核 CPU, 如服务器, 
%     如果是模拟出来的多线程, 即一个单核 CPU 的情况下, 在一个时间点只能执行一个代码, 
%     因为切换很快, 所以有同时执行的错觉. 
% \end{enumerate}
% \begin{remark}
%     使用多线程机制之后, \lstinline|main()| 方法结束之时主线程结束了, 其他线程还没有结束, 但没有主线程也不能运行, 
%     最起码, 现在的 Java 程序中至少有两个线程并发, 一个是垃圾回收线程, 一个是执行 \lstinline|main()| 方法的主线程. 
% \end{remark}

% % \subsection{多线程开发}

% % \subsubsection{并发程序设计}

% % \begin{enumerate}[label=(\arabic*)]
% %     \item 同一对象被多个线程同时操作: 
% %     \item 特点: 同时安排多个任务, 这些任务可以彼此穿插着进行: 
% %     有些任务可能是并行的, 比如去菜市场买菜、回家发邮件和去洗浴中心学外语的某些路是重叠的, 
% %     这是的确是在做三件事: 
% %     但进菜市场买菜、回家发邮件和进洗浴中心学外语三者是互斥的, 每个时刻只能完成其中一件. 
% %     换句话说, 并发允许两个任务彼此干扰. 
% % \end{enumerate}

% % \subsubsection{并行}

% % 你做你的事、我做我的事, 互不干扰同时进行. 

% % \subsubsection{串行}

% % 一个程序处理房钱进程, 按顺序接着处理下一个进程, 一个接着一个进行. 

% \subsection{线程的生命周期}

% Java 中线程的生命周期分为以下几个状态:
% \begin{enumerate}
%     \item 新建状态 (new): 当线程对象被创建但还没有调用其 \lstinline|start()| 方法时, 线程处于新建状态. 在这个阶段, 线程仅仅是一个 Java 对象, 没有被分配到实际的操作系统线程. 
%     \item 就绪状态 (Runnble): 线程一旦调用了 \lstinline|start()| 方法, 它就进入就绪状态. 在就绪状态下, 线程等待被调度器分配一个可运行的时间片, 以便运行它的 \lstinline|run()| 方法. 
%     \item 运行状态 (Runing): 当线程被调度并执行了其 \lstinline|run()| 方法时, 它进入运行状态. 在运行状态中, 线程执行其实际的工作. 
%     \item 阻塞状态 (Bloked): 线程可能会在某些情况下被阻塞, 例如等待某个锁、等待输入/输出等. 当线程处于阻塞状态时, 它暂时停止执行, 并等待特定的条件发生, 使得线程可以继续执行. 
%     \item 等待状态 (Waiing): 线程进入等待状态是因为调用了 \lstinline|Object.wait()|、\lstinline|Thread.join()| 或类似的方法, 使线程等待某个条件的发生. 
%     \item 超时等待状态 (Timed Waiting): 类似于等待状态, 但在指定的时间内会自动返回, 例如通过调用 \lstinline|Thread.sleep()|、\lstinline|Object.wait(long timeout)| 或 \lstinline|Thread.join(long millis)| 等方法. 
%     \item 终止状态 (Terminted): 当线程的 \lstinline|run()| 方法执行完成或者线程调用了 \lstinline|Thread.stop()| 方法时, 线程进入终止状态. 一旦线程终止, 它就不能再重新启动. 
% \end{enumerate}

% 线程状态之间的转换如下: 

% - 新建状态 -> 就绪状态: 调用 `start()` 方法. 
% - 就绪状态 -> 运行状态: 被调度器选中并执行 `run()` 方法. 
% - 运行状态 -> 阻塞状态、等待状态、超时等待状态: 调用了相关的方法. 
% - 阻塞状态、等待状态、超时等待状态 -> 就绪状态: 条件满足. 
% - 运行状态 -> 终止状态: `run()` 方法执行完成. 
% - 阻塞状态、等待状态、超时等待状态 -> 终止状态: 线程被中断或执行完成. 


% \subsection{线程同步}

% 当有一个线程在对内存进行操作时, 其他线程都不可以对这个内存地址进行操作, 直到该线程完成操作,
% 其他线程才能对该地址进行操作, 而其他线程又处于等待状态, 实现线程同步的方法有很多种.

% 为什么要线程同步? 多个线程同步运行的时候可能调用线程函数, 在多个线程同时对同一个内存地址进行
% 写入, 由于 CPU 时间调度上的问题, 写入数据会被多次的覆盖, 所以需要使线程同步, 防止数据并发
% 导致的安全性问题.

\chapter{继承和多态}

面向对象范式将数据和方法结合在对象中, 使用面向对象范式的软件设计聚焦于对象以及对象上的
操作. 面向对象方法结合了面向过程范式的强大之处, 并且进一步将数据和操作集成在对象中.

继承对于软件重用是一个重要且功能强大的特征.

\chapter{抽象类和接口}

\section{抽象类}

抽象方法不能包含在非抽象类中, 如果抽象父类的子类不能实现所有的抽象方法, 那么子类
也必须定义为抽象的. 在继承自抽象类的非抽象子类中, 必须实现所有的抽象方法. 此外, 抽象方法
是非静态的.

抽象类不能使用 \lstinline|new| 操作符来初始化, 但仍然可以定义它的构造方法, 这个构造方法
在其子类的构造方法中调用.

包含抽象方法的类必须是抽象的, 但可以定义一个不包含抽象方法的抽象类, 这个抽象类用作
定义新子类的基类.

子类可以重写父类的方法并将其定义为抽象的, 这是很少见的一种做法, 一般用于父类的方法实现
在子类中变得无效的情况. (尝试举例?)

即使父类是具体类, 子类也可以是抽象的.

不能使用 \lstinline|new| 操作符来从一个抽象类创建一个实例, 但是抽象类可以用作一种数据类型.

\section{接口}

接口的目的是指明相关或者不相关的类的对象的共同行为. 定义接口的语法为
\begin{lstlisting}[language=Java]
accessModifier interface InterfaceName {
    
}
\end{lstlisting}

在 Java 中, 接口可以看作是一种特殊的类, 你可以使用接口作为引用变量的数据
类型或者作为类型转换的结果类型, 但是不能使用 \lstinline|new| 操作符创建
接口的实例.

一个类可以使用 \lstinline|implements| 关键字来实现一个接口.

接口中修饰数据域的 \lstinline|public static final| 和修饰方法的
\lstinline|public abstract| 可以被忽略, 但是在实现该接口的类中,
必须将接口中的方法声明为 \lstinline|public| 的.

如果在声明接口方法时使用了关键字 \lstinline|default|, 那么你实际上
给出了一个默认的接口方法, 默认的接口方法需要在接口的定义中实现. 实现该接口
的类可以简单地使用默认的方法来实现, 或者重写该接口方法. 利用该特征,
在一个具有默认实现的已有接口中添加新的方法时, 只要声明默认方法并给出实现,
就可以不必为已经实现了该接口
的已有类重新编写接口方法的实现.

接口中的公共静态方法与类中的公共静态方法相同.

在实现接口的默认方法和公共静态方法时, 可以定义并调用接口的私有方法.

\begin{table}[H]
    \centering
    \caption{接口中的方法}
    \begin{tabular}{p{0.5\textwidth}p{0.45\textwidth}}
        \hline 
        \textbf{方法声明部分形如……} & \textbf{使用方法} \\ 
        \hline 
        \lstinline|public abstract returnType p()| & 等价于 \lstinline|returnType p()| \\
        \lstinline|public default returnType p()| & 默认接口方法, 实现该接口的子类可以选择直接使用它或者重写它 \\ 
        \lstinline|public static returnType p()| & 公共静态方法, 与类中的公共静态方法的使用方式相同 \\ 
        \lstinline|private returnType p()| & 私有方法, 可以被默认接口方法和公共静态接口方法调用 \\ 
        \hline
    \end{tabular}
\end{table}

\section{标记接口}

内部为空的接口称为\textbf{标记接口} (marker interface),
标记接口既不包括敞亮也不包括方法, 它用来表示一个类拥有某些希望具有的特征.

\section{抽象类和接口的区别与联系}

\begin{table}[H]
    \centering
    \caption{抽象类与接口}
    \begin{tabular}{cp{0.3\textwidth}p{0.25\textwidth}p{0.3\textwidth}}
        \hline 
        & \textbf{变量} & \textbf{构造方法} & \textbf{方法} \\
        \hline
        抽象类 & 无限制 & 子类通过构造方法链调用构造方法, 抽象类不能用 \lstinline|new|
        操作符实例化 & 无限制 \\ 
        接口 & 所有的变量必须是公开、静态、最终, 即 \lstinline|public static final| 的 &
        没有构造方法, 不能用 \lstinline|new| 操作符实例化 
        & 可以包含 \lstinline|public abstract| 的抽象实例方法、\lstinline|public default|
        的默认接口方法以及 \lstinline|public static| 的静态方法. 在实现默认接口方法
        与静态方法时, 作为辅助, 可以定义并调用 \lstinline|private| 的接口内部私有方法. \\ 
        \hline
    \end{tabular}
\end{table}

\part{数据结构与算法}

\chapter{线性表}

\section{链表}

链表是一种数据结构, 其中的元素 (节点) 通过指针相互连接. 本节我们介绍如何在 Java 中实现一个
简单的链表. 它主要有三个类构成 —— 实现节点的 \lstinline|Node| 类、
实现链表本身的 \lstinline|LinkedList| 类, 以及一个用于测试的编写了 \lstinline|main| 方法的测试类.

\subsection{\lstinline|Node| 类: 链表的节点}

% \lstinputlisting[language=Java]{../IdeaProjects/LearningJava/src/LinkedList/Node.java}

在类 \lstinline|Node| 中, 我们定义了两个属性: \lstinline|data| 属性表示我们存放在这个
节点中的数据的值, 而 \lstinline|next| 属性指的是当前节点的下一节点. 为了保证程序的封装性
不被破坏, 你可以把这些参数设置为私有 (\lstinline|private|), 但现在我们简单起见, 把这些
属性的访问修饰符都设置为缺省, 这意味着我们可以通过简单的赋值运算实现属性的访问和变更, 而不需要
使用公共的 \lstinline|getter| 和 \lstinline|setter|.

我们为 \lstinline|Node| 类提供了一个构造器, 它的作用是把节点的 \lstinline|data| 属性
设置为我们给定的形式参数 \lstinline|data|, 然后默认它是链表的尾节点, 即 \lstinline|this.next = null|.

\subsection{\lstinline|LinkedList| 类: 链表的实现}

我们先给出 \lstinline|LinkedList| 类的完整的实现:
% \lstinputlisting[language=Java]{../IdeaProjects/LearningJava/src/LinkedList/LinkedList.java}

\part{代码设计}

\chapter{代码设计概述}

学习代码设计的相关知识可以帮助我们写出可扩展、可读和可维护的高质量代码.

{\captionof{table}{代码质量的评判标准} % 表格标题
\label{代码质量的评判标准} % 交叉引用标签
\begin{longtable}{cp{0.4\textwidth}p{0.4\textwidth}}
    \toprule
    % 表头
    \textbf{特征} & \textbf{含义} & \textbf{评判标准} \\
    \midrule
    \endhead
    \bottomrule
    \endfoot

    % 表格内容
    可维护性 
    & 在不破坏原有代码设计、不引入新 bug 的情况下, 能够快速修改或添加代码 
    & 与多种因素有关: 代码简洁、可读性好、可扩展性好、分层清晰、模块化程度高、
    高内聚低耦合、遵守基于接口而非实现编程的设计原则 \\
    可读性
    & 代码是否易读、易理解
    & 代码符合代码规范、命名达意、注释详尽、函数长度合理、模块划分清晰、高内聚低耦合、Code Review \\ 
    可扩展性
    & 在不修改或少量修改原有代码的情况下, 能够通过扩展方式添加新功能
    & \\ 
    灵活性 & 含义比较宽泛 & \\ 
    简洁性 & 代码简单、逻辑清晰 & KISS 原则 \\ 
    可复用性 & 尽量减少重复代码的编写, 复用已有代码 & DRY 原则 \\ 
    可测试性 & 易于编写单元测试 & \\
\end{longtable}}

\section{关于面向对象编程范式的进一步讨论}

面向对象编程范式因其丰富的特性 (封装、抽象、继承和多态) 可以实现很多复杂的设计思路,
因此它是很多设计原则、设计模式编码实现的基础.

\subsection{面向对象编程范式的特性}

\subsubsection{封装 (encapsulation)}
 
封装也称为信息隐藏或数据访问保护. 类通过暴露有限的访问接口, 授权外部仅能通过类
提供的方式来访问内部信息或数据. 这看似影响了程序代码的灵活性, 但其实过度灵活意味着
不可控和不安全, 属性可以通过各种奇怪的方式被随意修改, 而且修改逻辑可能散落在代码的
各个角落, 进而影响代码的可读性和可维护性.

类通过提供有限的方法暴露必要的操作, 能提高类的易用性. 将属性封装, 暴露少许的必要的方法
给调用者, 调用者就不需要了解太多的业务细节, 用错的概率也会降低很多.

\subsubsection{继承 (inheritance)}

继承最大的作用是代码复用.

\subsubsection{多态 (polymorphism)}

多态是指, 在代码运行过程中我们可以用子类替换父类, 并调用子类的方法.

\newpage
\thispagestyle{empty}

\chapter{设计模式}

\section{创建型设计模式}

\subsection{单例模式}

% \begin{code}
%     \lstinline|Strik0r.java|
%     \lstset{
%         moreemph={INSTANCE},
%         emphstyle=\color{class_field},
%     }
%     \lstinputlisting[language=Java]
%     {../IdeaProjects/Hello/src/cn/edu/nwpu/software/strik0r/root/Strik0r.java}
% \end{code}
\part{编程速查}

\chapter{时间相关的需求……}

\section{\lstinline|Calendar| 类及其子类： 表示精确到毫秒的特定时刻}

先前我们曾经介绍过, 学习各种 Java 类的核心在于明确该类在对哪一种客观物质世界
的对象进行建模. 一个具体的 \lstinline|GregorianCalendar| 对象代表了一个具体
的公历 (gregorian calendar) 时间节点.

{\captionof{table}{\lstinline|Calendar| 类的域常量} % 表格标题
% \label{} % 交叉引用标签
\begin{longtable}{p{0.3\textwidth}p{0.5\textwidth}}
    \toprule
    % 表头
    \textbf{常量} & \textbf{说明} \\

    \toprule
    \endhead
    \bottomrule
    \endfoot

    % 表格内容
    \lstinline|YEAR| & 年份 \\ 
    \lstinline|MONTH| & 月份, 其中 $0$ 表示 $1$ 月 \\ 
    \lstinline|DATE| & 日前 \\ 
    \lstinline|OUR| & 小时 (12 小时制) \\ 
    \lstinline|HOUR_OF_THE_DAY| & 小时 (24 小时制) \\ 
    \lstinline|MINUTE| & 分钟 \\ 
    \lstinline|SECOND| & 秒 \\ 
    \lstinline|DAY_OF_WEEK| & 一周中的哪一天, 1 表示星期天 \\ 
    \lstinline|DAY_OF_MONTH| & 同 \lstinline|DATE| \\ 
    \lstinline|DAY_OF_YEAR| & 一年中的哪一天, 其中 1 表示该年的第一天 \\ 
    \lstinline|WEEK_OF_MONTH| & 一月中的哪一周, 其中 1 表示该月的第一周 \\ 
    \lstinline|WEEK_OF_YEAR| & 一年中的哪一周, 其中 1 表示该年的第一周 \\ 
    \lstinline|AM_PM| & $0$ 表示上午, $1$ 表示下午 \\ 
\end{longtable}}

\subsection{复习参考题}

\begin{exercise}
    可以使用 \lstinline|Calendar| 类来创建一个 \lstinline|Calendar| 对象吗?
    这显然是不可以的, 因为 \lstinline|Calendar| 类是一个抽象类, 不能够
    通过 \lstinline|new| 操作符构造一个 \lstinline|Calendar| 对象.
\end{exercise}

\appendix
% 设置 chapter 标题样式
\titleformat{\chapter}[hang]{\centering\heiti\Large\bfseries}{附录\,\thechapter}{1em}{}
\renewcommand{\chaptermark}[1]{\markboth{附录 \thechapter\, #1}{}}

\newpage
\begin{thebibliography}{1}
    \addcontentsline{toc}{chapter}{参考文献}
    \bibitem{核心技术1}
    [美] Cay S. Horstmann. 
    Java 核心技术: 原书第 12 版. 卷 I: 开发基础[M].
    林琪, 苏钰涵译.
    北京: 机械工业出版社, 2022
    \bibitem{核心技术2}
    [美] Cay S. Horstmann.
    Java 核心技术: 原书第 12 版. 卷 II: 高级特性[M].
    陈昊鹏译.
    北京: 机械工业出版社, 2023
    \bibitem{设计模式之美}
    王争. 设计模式之美 [M]. 北京: 人民邮电出版社, 2022.
\end{thebibliography}

\newpage
\thispagestyle{empty}
\vspace*{5cm}
\begin{center}
    \includegraphics*[width=\textwidth]{pic/i_love_npu.jpeg}
    \large
    公诚勇毅 \quad 永矢毋忘

    中华灿烂 \quad 工大无疆
\end{center}
\vspace*{13em}
\begin{center}
    \small
    本文档由\textbf{钱锋}编写, 钱锋保留一切权利.

    文档中出现的部分素材来源于网络, 笔者承诺这些素材仅供学习交流之用, 
    它们的原作者保留一切权利.

    2023 年 \quad 西北工业大学 \quad 中国西安 
\end{center}

\end{document}